\newpage
\chapter{Automating/Scripting the model}
There are three main ways to automate the mode. The first is the parameter scan window, see section \ref{sec:scanwindow}.  This alows the user vary a paramter in steps. The second way is by using Python scripting, see section \ref{sec:pythonscripts} and the third way is through matlab scripting see section \ref{sec:matlabscripts}. 

\section{The parameter scan window}
\label{sec:scanwindow}

Sometimes one wishes to systematically vary a simulation parameter, to do this first bring up the parameter scan window, this can be done by clicking on the 




\begin{figure}[H]
\centering
\includegraphics[width=\textwidth]{./images/param_scan.png}
{\caption{Step 1: Select the 'Parameter scan' tool.}}
\label{overflow}
\end{figure}

Then make a new scan by clicking on the new button (1) then open the new scan by double clicking on the icon representing the scan (2). See figure \ref{fig:newscan}. This will bring up the scan window, see figure \ref{fig:newscanline}.

\begin{figure}[H]
\centering
\includegraphics[width=0.5\textwidth]{./images/param_scan_new.png}
\caption{Step 2: Make a new parameter scan, then double click on it to open it.}
\label{fig:newscan}
\end{figure}

\subsection{Changing one material parameter}

Once the scan window has appeared. Make a new scan line by clicking on the the plus icon \ref{fig:newscanline} (1), then select this line so that it is highlighted (2), then click on the three dots (3) to select which parameter you want to scan. In this example we will be selecting the electron mobility of a P3HT:PCBM solar cell. Do this by navigating to epitaxy$\rightarrow$ P3HT:PCBM$\rightarrow$ Drift diffusion$\rightarrow$ Electron mobility y. Highlight the parameter and then click OK. This should then appear in the scan line. The meaning of "epitaxy$\rightarrow$ P3HT:PCBM$\rightarrow$ Drift diffusion>Electron mobility y" will now be explained below:

\begin{itemize}
  \item epitaxy: All parameters in the .gpvdm file are exposed via the parameter selection window see \ref{fig:scanselect}. This file is a tree structure, see \ref{sec:gpvdmfileformat}. The all parameters which define the device it's self are contained under epitaxy.
  \item P3HT:PCBM: Under epitaxy each layer of the device is given by its name. The active layer in this device is called P3HT:PCBM.
  \item Drift diffusion: All electrical parameters are stored under drift diffusion.
  \item Electron mobility y: One can define asymmetric mobilities in the z,x and y direction - this is useful for OFET simulations.  However by default the model assumes a symmetric mobility which is the same in all directions. This value is defined by "Electron mobility y". 
\end{itemize}

\begin{figure}[H]
\centering
\includegraphics[width=\textwidth]{./images/param_scan_new_line.png}
\caption{Step 3: Add a 'scan line' to the scan.}
\label{fig:newscanline}
\end{figure}

\begin{figure}[H]
\centering
\includegraphics[width=0.5\textwidth]{./images/param_scan_select.png}
\caption{Step 5: Select the parameter you want to scan in the parameter selection window, in this case we are selecting epitaxy$\rightarrow$ P3HT:PCBM$\rightarrow$ Drift diffusion$\rightarrow$ Electron mobility y.}
\label{fig:scanselect}
\end{figure}

Next enter the values of mobility which you want to scan over in this case we will be entering "1e-5 1-6 1e-7 1e-8 1e-9" (see figure \ref{fig:runscan} 1) then click "run scan" (see figure \ref{fig:runscan} 2). Gpvdm will run one simulation on each core of your computer until all the simulations are finished.

\begin{figure}[H]
\centering
\includegraphics[width=\textwidth]{./images/param_scan_inputvalues.png}
\caption{Step 6: Enter the input values of mobility (or other values) you want to scan over (1). Then run the simulations.}
\label{fig:runscan}
\end{figure}

To view the simulation results click on the "output" tab this will bring up the simulation outputs, see figure \ref{fig:scanoutput}. You can see that a directory has been created for each variable that we scanned over so 1e-5, 1-6, 1e-7, 1e-8 and 1e-9.  If you look inside each directory it will be an exact copy of the base simulation directory.  If you double click on the files with multi colored JV curves, see the red box in figure \ref{fig:scanoutput}. Gpvdm will automaticity plot all the curves from each simulation in one graph, see figure \ref{fig:scanjv}.

\begin{figure}[H]
\centering
\includegraphics[width=\textwidth]{./images/param_scan_output.png}
\caption{Step 7: The output tab.}
\label{fig:scanoutput}
\end{figure}

\begin{figure}[H]
\centering
\includegraphics[width=0.5\textwidth]{./images/param_scan_jv.png}
\caption{Step 8: The result of the mobility scan.}
\label{fig:scanjv}
\end{figure}

\subsection{Duplicating parameters - changing the thickness of the active layer}

Very often one wants to change a parameter, then set another parameter equal to the parameter which was changed. An example of this is one may want to change electron and hole mobilities together when simulating a device with symmetric mobilities. This can be done using the duplicate function of the scan window as seen in figure \ref{fig:scanduplicate}.  In this example we tackle a slightly more trick problem than changing mobilities together we are going to change the physical width of the active layer and at the same time adjust the electrical mesh to make it match.  As discussed in section \ref{ref:mesh} the width of the active layer must always match the width of the electrical mesh.  When you change the layer width by hand in the layer editor gvpdm updates the width of the electrical mesh for you. But when scripting the model it won't do this update for you.  Therefore in the example below we are going to set the width of the active layer by scanning over:

epitaxy$\rightarrow$P3HT:PCBM$\rightarrow$dy of the object
\\
\\
Then we are going to add another line under and under parameter to scan select
\\
\\
mesh$\rightarrow$mesh\_y$\rightarrow$segment0$\rightarrow$len
\\
\\
and set it to
\\
\\
epitaxy$\rightarrow$P3HT:PCBM$\rightarrow$dy of the object
\\
\\
under the operation dropdown box. You will see the word duplicate appear under values.
\\
\\
If you now run the simulation "epitaxy$\rightarrow$P3HT:PCBM$\rightarrow$dy of the object" will be changed and "mesh$\rightarrow$mesh\_y$\rightarrow$segment0$\rightarrow$len" will follow it.


\begin{figure}[H]
\centering
\includegraphics[width=0.5\textwidth]{./images/param_scan_duplicate.png}
\caption{Duplicating material paramters.}
\label{fig:scanduplicate}
\end{figure}

\subsection{Setting constants}
Often when running a parameter scan one wants to set a constant value, this can be done using the "constant" option in the Operations dropdown menu. See figure \ref{fig:scanconst}

\begin{figure}[H]
\centering
\includegraphics[width=0.5\textwidth]{./images/param_scan_const.png}
\caption{The result of the mobility scan.}
\label{fig:scanconst}
\end{figure}

\subsection{The equivalent of loops}
Often when scanning over a parameter range one may want to simulate so many parameters that it is not practical to type them in.  In this case gpvdm has the equivalent of a loop. So for example if one wanted to change a value from 100 to 400 in steps of 1, one could type

\begin{listing}[H]
\begin{minted}[frame=single,
               framesep=3mm,
               linenos=false,
               xleftmargin=21pt,
               tabsize=4]{matlab}

[100 400 1]

\end{minted}
\caption{JSON example} 
\label{json-example}
\end{listing}

\subsection{Limitations of the scan window}
Although the scan window is convenient in that it provides a quick way to scan simulation parameters, it is by nature rather limited in terms of flexibility. If you want to do complex scans were multiple parameters are changed or to programmatically collect data from each simulation then you can use the \index{python} or matlab interfaces to gpvdm.  These are described in the next section.

\section{Python/MATLAB scripting of gpvdm}
Scripting offers a more powerful way to interact with gvpdm. Rather than using the graphical user interface, you can use your favourite programming language to interact with gpvdm.  This gives you the option to drive gpvdm in a far more powerful way than can be done using the graphical interface alone.  Below I give examples of using MATLAB and python to drive gpvdm, but you can use any language you want which has a json reader/writer.  Pearl and Java are two languages which spring to mind.

Before you begin scripting gpvdm, open up the install path of gpvdm, the default gpvdm will be installed to C:\textbackslash Program files x86 \textbackslash gpvdm, in there you will see in this directory there are two windows executables, one called \emph{gpvdm.exe}, this is the graphical user interface, and a second .exe, called \emph{gpvdm\_core.exe}.  You can run \emph{gpvdm\_core.exe} from the command line without \emph{gpvdm.exe}. You simply need to navigate to a directory containing a \emph{sim.gpvdm} folder and call \emph{gpvdm\_core.exe}, this can be done from the windows command line, matlab, python or any other scripting language.
However, before you can do this on windows, you need to add C:\textbackslash Program files x86 \textbackslash gpvdm to your windows path so that windows knows where gpvdm is instaled.  An example of how to do this on a modern version of windows is given in the link
\url{https://docs.microsoft.com/en-us/previous-versions/office/developer/sharepoint-2010/ee537574(v=office.14)}

Every new version of windows seems to move the configuration options around, so you may have to find instructions for your version of windows.

\subsection{Python scripting}
\label{sec:pythonscripts}
Related YouTube videos:
\begin{figure}

\begin{tabular}{ c l }

\includegraphics[width=0.05\textwidth]{./images/youtube.png}

&
\href{https://www.youtube.com/watch?v=vyeAzxBZjMg}{Python scripting perovskite solar cell simulation}

\end{tabular}
\end{figure}

There are two ways to interact with .gpvdm files via python, using native python commands or by using the gpvdm class structures, examples of both are given below.


\subsubsection{The native python way}
As described in section \ref{sec:gpvdmfileformat}, .gpvdm files are simply json files zipped up in an archive. If you extract the sim.json file form the sim.gpvdm file you can use Python's json reading/writing code to edit the .json config file directly, this is a quick and dirty approach which will work. You can then use the $os.system$ call to run \emph{gpvdm\_core.exe} to execute gpvdm.

For example were one to want to change the mobility of the 1st device layer to 1.0 and then run a simulation you would use the code listed in listing \ref{python-example}.

\begin{listing}
\begin{minted}[frame=single,framesep=3mm,linenos=false,xleftmargin=21pt,tabsize=4]{python}

	import json
	import os
	import sys

	f=open('sim.json')		#open the sim.json file
	lines=f.readlines()
	f.close()
	lines="".join(lines)	#convert the text to a python json object
	data = json.loads(lines)

	#Edit a value (use firefox as a json viewer
	# to help you figure out which value to edit)
	# this time we are editing the mobility of layer 1
	data['epitaxy']['layer1']['shape_dos']['mue_y']=1.0


	#convert the json object back to a string
	jstr = json.dumps(data, sort_keys=False, indent='\t')

	#write it back to disk
	f=open('sim.json',"w")
	f.write(jstr)
	f.close()

	#run the simulation using gpvdm_core
	os.system("gpvdm_core.exe")
\end{minted}
\caption{Manipulating a sim.json file with python and running a gpvdm simulation.} 
\label{python-example}
\end{listing}

If the simulation in sim.json is setup to run a JV curve, then a file called sim\_data.dat will be written to the simulation directory containing paramters such as PCE, fill factor, $J_{sc}$ and $V_{oc}$.  This again is a raw json file, to read this file in using python and write out the value of $V_oc$ to a second file use the code given in listing \ref{python-example2}.

\begin{listing}
\begin{minted}[frame=single,framesep=3mm,linenos=false,xleftmargin=21pt,tabsize=4]{python}
f=open('sim_info.dat')
lines=f.readlines()
f.close()
lines="".join(lines)
data = json.loads(lines)

f=open('out.dat',"a")
f.write(str(data["Voc"])+"\n");
f.close()

\end{minted}
\caption{Reading in a sim\_data.dat file using Python's native json reader.} 
\label{python-example2}
\end{listing}



\subsubsection{Using gpvdm's built in classes for reading and writing json}
Gpvdm has a set of classes that can read in gpvdm files and write them to disk. The difference between using python's native commands and the gpvmd classes is that, gpvdm will convert the json save files to a hierarchical tree of python classes rather than leaving them as raw json. So for example using Python's native json interpreters one would write:

 
\begin{listing}
\begin{minted}[frame=single,framesep=3mm,linenos=false,xleftmargin=21pt,tabsize=4]{python}
	data['epitaxy']['layer1']['shape_dos']['mue_y']=1.0
\end{minted}
\caption{Reading in a sim\_data.dat file using Python's native json reader.} 
\label{python-example3}
\end{listing}

but using the gpvdm interpreter one would write

\begin{listing}
\begin{minted}[frame=single,framesep=3mm,linenos=false,xleftmargin=21pt,tabsize=4]{python}
	data.epitaxy.layer[1].shape_dos.mue_y=1.0
\end{minted}
\caption{Reading in a sim\_data.dat file using Python's native json reader.} 
\label{python-example4}
\end{listing}

The gpvdm class tree also has embedded functions for searching for objects and alike some of which are described below in listing \ref{python-example5}.

\begin{listing}
\begin{minted}[frame=single,framesep=3mm,linenos=false,xleftmargin=21pt,tabsize=4]{python}
#!/usr/bin/env python3
import json
import os
import sys

sys.path.append('c:\Program files x86\gpvdm\modules')
from gpvdm_json import gpvdm_data

data=gpvdm_data()
data.load("sim.json")
data.epitaxy.layer[1].shape_dos.mue_y=1.0
data.save()

os.system("gpvdm_core.exe")

\end{minted}
\caption{Editing sim.json files using gpvdm's built in classes.} 
\label{python-example5}
\end{listing}

\subsubsection{Running gpvdm across multiple cores}
The scan window by default uses gpvdm's built in job scheduler so that if you want to scan across 10 parameters and have a CPU with multiple cores, the jobs will be spread across all cores.  This increases the overall speed of the simulations. You can access this API using the gpvdm\_api class, an example of how to do this is given in listing \ref{python-example6}.

\begin{listing}
\begin{minted}[frame=single,framesep=3mm,linenos=false,xleftmargin=4pt,tabsize=4]{python}
#!/usr/bin/env python3
import os
import sys
sys.path.append('c:\Program files x86\gpvdm\modules')

from gpvdm_api import gpvdm_api

#initialize the API
api=gpvdm_api(verbose=False)

#Use the name of the current script to determine the directory name to make
script_name=os.path.basename(__file__).split(".")[0]

#define the name of the simulation dir
scan_dir=os.path.join(os.getcwd(),script_name)

#make the simulation dir
api.mkdir(scan_dir)					#make a new scandir

#tell the API where we are going to run the simulation
api.server.server_base_init(scan_dir)

#Loop over electron and hole mobilities.
for mue in [ 1e-5, 1e-6, 1e-7, 1e-8]:
	for muh in [ 1e-5, 1e-6, 1e-7, 1e-8 ]:
		#define the sub sim path
		sim_path=os.path.join(scan_dir,"{:.2e}".format(mue),"{:.2e}".format(muh))
		#make the directory
		api.mkdir(sim_path)

		#clone the current sim dir to the new dir
		api.clone(sim_path,os.getcwd())		

		#make edit the newly generated sim.json file
		data=gpvdm_data()
		data.load(os.path.join(sim_path,"sim.json"))
		data.epitaxy.layer[1].shape_dos.mue_y=mue
		data.epitaxy.layer[1].shape_dos.muh_y=muh
		data.save()

		#Add the path to the job list
		api.add_job(path=sim_path)

#run all the jobs over multiple CPUs
api.server.simple_run()

#Generate GNUPLOT compatible files for plotting the results together.
api.build_multiplot(scan_dir,gnuplot=True])

\end{minted}
\caption{Running jobs across multiple CPUs using python} 
\label{python-example6}
\end{listing}

\newpage
\subsection{MATLAB scripting}
\label{sec:matlabscripts}
As described in section \ref{sec:gpvdmfileformat} gpvdm simulations are stored in .json files zipped up inside a zip archive. Matlab has both a zip decompressor and a json decoder.  Therefore it is straight forward to edit and read and edit .gpvdm files in MATLAB. You can then use MATLAB to perform quite complex parameter scans.  The example script below in listing \ref{matlab-example} demonstrates how to run multiple simulations with mobilities ranging from 1e-7 to 1e-5 $m^{2}V^{-1}s^{-1})$. The script starts off by unzipping the sim.json file, if you already have extracted your sim.json file from the sim.gpvdm file you don't need these lines. The code then reads in sim.json using the MATLAB json decoder $jsondecode$. A new directory is made which corresponds to the mobility value, the sim.gpvdm file copied into that directory. Then $json\_data.epitaxy.layer0.shape\_dos.mue_y$ is set to the desired value of mobility and the simulation saved using $jsonencode$ and $fopen,fprintf,fclose$.  The $system$ call is then used to run $gpvdm\_core.exe$ to perform the simulation. Out put parameters such as $J_{sc}$ are stored in sim\_data.dat again in json format, see section \ref{sec:siminfo}, although this is not done in this simple script.



\begin{listing}
\begin{minted}[frame=single,
               framesep=3mm,
               linenos=false,
               xleftmargin=21pt,
               tabsize=4]{matlab}


if exist("sim.gpvmd", 'file')==false
 sprintf("No sim.gpvdm file found"); %Check if we have a sim.gpvdm file
end

if exist("sim.json", 'file')==false
 unzip("sim.gpvdm")  %if we don't have a sim.json file
					 %try to extract it
end

A = fileread("sim.json");  %Read the json file.
json_data=jsondecode(A);  %Decode the json file

mobility=1e-7  %Start mobility
origonal_path=pwd  %working with the current dir
base_dir="mobility_scan"  %output directory name
while(mobility<1e-5)
    dir_name=sprintf("%e",mobility);
    full_path=fullfile(origonal_path,base_dir,dir_name)	%join paths
    mkdir(full_path)  %make the dir
    cd(full_path)  %cd to the dir

	%Update the json mobility
    json_data.epitaxy.layer0.shape_dos.mue_y=mobility  %Change mobility
													   %of layer0
    
    copyfile(fullfile(origonal_path,"sim.gpvdm"),\\
		fullfile(origonal_path,base_dir,dir_name,"sim.gpvdm"))

	%now write the json file back to disk
    out=jsonencode(json_data);
    json_data
    fid = fopen("sim.json",'w');
    fprintf(fid, '%s', out);
    fclose(fid);

	%run gpvdm - This won't work if you have not added the gpvdm
	%install directory to your windows paths 
    system("gpvdm_core.exe")
    
	%Multiply mobility by 10
    mobility=mobility*10;
end

%Move back to the original dir
cd(origonal_path)

\end{minted}
\caption{An example of how to call gpvdm from MATLAB} 
\label{matlab-example}
\end{listing}

